% Generated by D:/tools/scratch_qdd/mmrec/render.py from data_c.py - edit the data, not this file.
\subsection{C. 의도 공유 프로토콜과 결합 --- 상위는 구간과 할 일을, VLA는 그 안의 `언제'를}\label{rec:C}
\noindent\textbf{한 줄:} VLA를 흔한 VLA처럼 쓰지 않고 말단 보정기·세부 조종 스틱으로 두는 틀(§87·§90·§91·§92)을 세우고 결합 배선(Task 1--26)을 끝냈다. 무료 사전 실행에서 VLA 단독이 0/12였고 그 중 8편은 런타임 교착이었다(수정함). 상위$\leftrightarrow$VLA 의도 공유는 필수(§96 보충 2). 남은 것은 학습 --- ser-A-min-3 VLA 재학습, 결합 형식 상위 학습, 결합 폐루프 검증.

\paragraph{C1. 틀: VLA를 VLA처럼 쓰지 않는다}
{\small 09-25 23:18 -- 09-26 01:07 KST}
\begin{itemize}
  \item \textbf{사용자} (ul 75, 09-25 23:18 KST): \intent{그리고 한가지 궁금한게 이게 VLA를 따르기는 하는데 VLA는 일반적으로 VLM모델을 엔드투 엔드학ㄱ습을할때 이 엔드포인트를 어디로 움직여야한다는 모델로서의 기능이 아니라 테스크 전체에 활용하는 vlm이잖아? 근데 우리는 이 vlm의 역할을 단순히 오락실 게임기처럼 위아래 앞뒤와 같은 것중에 선택을 하게 하는것에 활용을 한다는 점임을 강ㅇ조를 해야할것 같고 실제로도 그렇게 활용될 수 있도록 하고 있는거지?}
  \item \textbf{사용자} (ul 77, 23:50 KST): \intent{내가 기존의 VLA 모델에서 문제라고 봤던점은 지금 내가 하는 행동 의 단위 자체가 모든 동작들마다 일관적일 수 없다는 점이야. 생각을 해보자. 우리가 옷을접을때 그냥 책상 위에 있을때는 잘 접을 거야. 그런데, 책상위에 벽돌 30개가 있어서 옷을 잡으러 갈때 계속 방해해버려 VLA는 이걸 절대 못풀어. 그이유는 VLA는 행동의 연속성을 의미단위로 풀려하지만 그 행동들은 의미를 가지기에는 너무 짧은 구간들만 보고 있는거거든 볼거면 전체의 흐름을 모두 본다음에 그 의미와 행동을 매칭해야하는데 그게 아니잖아. 즉 이행동을1}\textasciitilde\intent{3초간의 이미지만 가지고 내가 지금 뭐하고 있는거게}\textasciitilde\intent{ 하는 거나 다름이 없는 이상한 거야. 그래서 효율적인 학습을 위해선 이 전체적인 학습의 흐름이 뭐다 라는걸 아스트라가 잡아주는 과정이 굉장히 중요하다고 생각하는거야. 이걸 어떻게 유기적으로 활용할 수 있을까? 나랑 지금 브레인 스토밍해보자}
  \item \textbf{사용자} (ul 77, 23:50 KST): \intent{내가 방해물을 예시로 들었지만 그게 아니라 그냥 다른환경에 취약하다는 얘기야}
  \item \textbf{결과:} 가설 보정 조사(a8d4242, 09-26 00:10): 좁은 시연으로 순진하게 끝단 학습하면 국소 지름길을 배워 새 위치·시점에서 무너지는 경향은 체제 한정 지지; `무조건 과적합'은 반박; `짧은 창이 원인'은 근거 없음. 조향 표현 조사(1079716): 동작 수준 명령 인터페이스는 지지, `어디로' 공간 목표는 없음.
  \item \textbf{사용자} (ul 81--82, 09-26 00:46--00:50 KST): \intent{아게 진짜로 잘 모르긴하겠다 아스트라를 어케 쓰지?}
  \item \textbf{사용자} (ul 81--82, 09-26 00:46--00:50 KST): \intent{나는 VLA를 VLA 처럼 쓰고 싶지 않은거야 뭔소리인지 알아? 단점이 있기에 그건 순간 순간 어디로 가야하는지 에ㅐㄴㄴ드포인트 정밀 보정하는 역할을 하는 용도로 쓸거여}
  \item \textbf{사용자} (ul 81--82, 09-26 00:46--00:50 KST): \intent{그래서 전체적으로 다 엔드포인트를 위한 움직임 만을 하난 파이프 라인인거야. 이해했어? 근데 그럼 계획을 VLA로 부터 뺐어 왔으니까 이걸 아스트라가 담당해야하는데 아스트라가 잘하는 걸 하게 해야한다는거지}
  \item \textbf{그래서(결정):} 논문 중심 틀(cd5210f, 01:07 논문 갱신): VLA = 순간순간의 말단 정밀 보정기, VLM은 오락실 조이스틱처럼 typed 보기만 고름; 과제 이해·흐름·계획은 위에서(Astra J1 + 흐름 교정, A1·A2).
\end{itemize}

\paragraph{C2. 결합 설계 승인, MolmoAct를 주 참고로, 구현 계획}
{\small 09-26 01:44 -- 03:00 KST}
\begin{itemize}
  \item \textbf{사용자} (ul 84, 01:44 KST): \intent{몰모 엑트를 많이 참조분해해서 잘사용 적용해보면 좋을듯}
  \item \textbf{사용자} (ul 85, 02:11--02:13 KST): \intent{1}\textasciitilde\intent{3 다 허용할게 몰모엑트 많이 참고하고}
  \item \textbf{사용자} (ul 85, 02:11--02:13 KST): \intent{1,2,3 허용했지만 옳은 방향으로 바꿔도돼 확인해보고}
  \item \textbf{그래서(결정):} 정본 §84(c915535, 02:11): 결합 설계 승인, E-MA1 승인, 같은 진단 반복 시 falsify 먼저 평가, MolmoAct 채택 순위(미래 손끝 궤적 보조 $\rightarrow$ 이력 덧그림·궤적 조건 $\rightarrow$ 층별 KV $\rightarrow$ 깊이 보조 보류). 보충 1(a5816e1, 02:13): falsify 재사용은 같은 에피소드·같은 실패 서명에서 한 번까지, E-MA1 해석 주의(일반화 근거로 쓰지 않음), 호출 빈도는 비용 한도로.
  \item \textbf{그래서(결정):} 구현 계획(bb1f5f7, 02:59)과 메인 판정(7bb9035, 03:00): 두 층 합의 + T1 전제, 탐침 예산 15,000원, 실행은 R7 E2E 준비 기준점 뒤 subagent-driven.
  \item \textbf{다음:} R7 E2E 준비 기준점 도달(36·37회차 연속 무결, 태그 \texttt{stage3-\allowbreak{}e2e-\allowbreak{}ready} = 3d4738a, 8de6b80 09:46) 뒤 구현 착수(C5).
\end{itemize}

\paragraph{C3. 잡기·놓기 판단은 VLA의 선택지, VLA 자체가 세부 조종 스틱}
{\small 09-26 11:14 -- 11:17 KST (구현 15:17)}
\begin{itemize}
  \item \textbf{사용자} (ul 93, 11:14 KST): \intent{근데 그냥 소능 움직일때 물체가 근접하지 않았다면 상관 없지만 물체에 가깝게 다가가서 미세조정이 필요한데 아스트라의 반응성이 너무 느리면 그걸 보조하는 역할이라서 어떤걸 잡아야한다 이때 놓아야한다등에 대한 지식을 가져야하는데}
  \item \textbf{사용자} (ul 94, 11:17 KST): \intent{그 잡고, 놓고 하는 판단 자체도 vla가 선체택지에 있어야한다는 뜻}
  \item \textbf{사용자} (ul 94, 11:17 KST): \intent{vla자테가세부 조종스틱이기도하지 물론}
  \item \textbf{그래서(결정):} 거리에 따른 권한 분배(a4ae36e): 먼 이동 구간은 조이스틱·Astra를 강하게 따르고, 접촉 근처(\texttt{core.\allowbreak{}near\_\allowbreak{}contact} $\leq$ 5 cm)는 VLA 자기 과제 지식이 주도. 정본 §87(6e9f778): 융합 VLA typed 결정에 그리퍼 질문(쥐기 close / 놓기 open / 유지 keep); 선택지는 `의도', 실행 허가는 두 층 관문 + T1. 계층(6127f92): Astra = 느린 큰 조종, VLA 결정 = 세부 조종 스틱(0.33 s), 행동 전문가 = 그 입력을 연속 움직임으로.
  \item \textbf{결과:} 구현 확인(Task 12, 340c5dc 15:17): 그리퍼 질문은 어디에도 없었고 \texttt{phase}는 쥐기·놓기 의도를 명시하지 않았다 $\rightarrow$ 그리퍼 질문을 융합 VLA에만 추가, 라벨 규칙은 \texttt{harvest/\allowbreak{}intent.\allowbreak{}py} 한 곳(R2 = 기록 단계, S-E2E = 그리퍼 명령 속도 0.176 /s).
\end{itemize}

\paragraph{C4. 구간 계획(§90) · 도착 시 대조(§91) · 구조 확정(§92)}
{\small 09-26 14:37 -- 14:48 KST}
\begin{itemize}
  \item \textbf{사용자} (ul 100, 14:37 KST): \intent{상위 아스트라가 어디로 어케 움직일거고 나는 지금 다가가는 행동ㅈ이다, 이 구간선 잡는걸 할거다. 이구간선 드는걸할거다 하는ㅇ계회도 담고 vla는 그 구간이 길기에 그구간에서 잡는게 있다곤 하는데 언제 잡을지를 상황을보고 결정하는 겆ㄱ}
  \item \textbf{그래서(결정):} 정본 §90(aed65ff, 14:37): Astra 답에 구간 계획(지금 구간·이 구간에서 할 일·다음 구간), VLA는 구간 의도를 입력으로 받아 그 안에서 `언제'를 스스로 정함(공간 목표는 여전히 없음).
  \item \textbf{사용자} (ul 101, 14:40 KST): \intent{그리고 아스트라의 판단이 몇초전 이미지와 상황을 보고한거지 그리고 실제vla가 조정한건 어땠는지의 비교를가지고도 추가판단이 가능해야함}
  \item \textbf{그래서(결정):} 정본 §91(5158cbc, 14:40): 도착 시 대조 --- 판단 기준 시각의 상태와 지금, Astra 명령과 그사이 VLA의 실제 움직임을 비교해 이미 됨 / 아직 유효 / 상황 바뀜 / 충돌을 판정, 다음 요청 문맥과 추가 학습 기록에 넣음.
  \item \textbf{사용자} (ul 102, 14:44--14:48 KST): \intent{이제 어떻게 서로 엮어서 최대 효율을 내게끔 파이프라인 구조 확정해서 이미지좀 만들어줘봐봐}
  \item \textbf{사용자} (ul 102, 14:44--14:48 KST): \intent{아직 돌고 있는 작업들은 확정이 나야지 파이프라인 확정이지? 아직 확실하지 않은 전제들이 좀 있지?}
  \item \textbf{그래서(결정):} 정본 §92(c410ec2, 14:44): ① Astra(low) 직렬 흐름 ② 도착 시 대조 + 합의 + 권한 a ③ VLA 결정 = 세부 조종 스틱 ④ 행동 전문가 ⑤ 안전 관문 ⑥ 실행 ⑦ 기록 + 그림. 보충 1(553a8ad, 14:48): `배치' 확정일 뿐, 전제 P1--P12(뒤에 P13 H-L)가 검증되기 전까지 잠정.
  \item \textbf{참고:} 뒤 결과로 본 전제(정본 표의 상태 칸은 작성 당시 값): P1 실데이터 준수 0.8 --- E-SR1d 0.536으로 불성립, E-SR1e 진행 중(D7·D8); P3 결합 > VLA 단독 --- VLA 단독 0/12로 짝 비교 보류(§92 보충 2, C7); P4 명령 정확도 --- 수정 방향 13/13, 7--10 cm 알아채기 0/13(A6·A7); P10 E-Astra-necessity --- 초안, 유료 금지; P13 H-L --- 8B 경향 4/4(B6), 35B 미정.
\end{itemize}

\paragraph{C5. 결합 배선 구현 Task 1--26과 리뷰에서 잡은 결함}
{\small 09-26 10:02 -- 23:59 KST}
\begin{itemize}
  \item \textbf{해 봄:} 계획 \texttt{docs/\allowbreak{}superpowers/\allowbreak{}plans/\allowbreak{}2026-\allowbreak{}09-\allowbreak{}26-\allowbreak{}astra-\allowbreak{}vla-\allowbreak{}coupling.\allowbreak{}md}, subagent-driven, 통제자 판정. Task 1 파라미터·답 스키마·\texttt{astra-\allowbreak{}couple@v1}(205db33 10:02), 2 답 게이트(7c56331), 3 가격표·공유 장부 80 \% 정지(0e480ac), 4 직렬 흐름·모의 Astra(e6f163c), 5 두 답 이력(e446070), 6 부드러운 편향(b5c059b), 7 두 층 비가역 관문(T1 근거만)·VLA 빠른 검증(f9eb416), 8 Astra 전용 덧그림·카메라 모델(8350c7e·85bd2ed), 9 CoupleDriver(dd3ce0c), 10 OursRuntime 배선(dd1afbc), 11 falsify 먼저 재사용(8074683), 12 직렬화 \texttt{ser-\allowbreak{}A-\allowbreak{}min-\allowbreak{}3}(340c5dc 외), 13 로컬 VLM 어댑터(58eed4b), 14 폐루프 결합 팔·유료 가드(ea98034), 15·16 E-Couple·E-Astra-necessity 등록 초안(2a99db6·2eea17b), 17 편향 $\times$ 권한 a(b0ee916), 18 \texttt{astra-\allowbreak{}couple@v2}(8c9064a), 19 도착 시 대조(20d735c), 20 프롬프트 검진 가드(a913020), 21 사건 불응기·미응답 보고·quota 치명 정지(688992c), 22 확신도 기록 훅(5a940d2), 25 교착 수정(c835476), 26 유료 가드 ul 114(99b3b58), 최종 검토 수정(fddefd4 23:59).
  \item \textbf{결과:} 리뷰에서 잡아 고친 결함(통제 판정): 구간 줄이 그리퍼 정답을 알려 주던 것(F12, bce0bd8 --- 구간이 사건을 품게, 학습 때만 unknown 구간 p 0.3); 권한 a가 떨어질 때 남은 편향이 1.92 cm·0.36 rad 튀던 것(F17, badb1e2); 실행 청크가 없을 때 결정 중심으로 거짓 화살표를 그리던 것(F18, 9328064); 늦은 답이 반영되던 것·모듈형 스킬이 편향을 매 틱 지워 가짜 충돌을 내던 것(F19·F19b·F19c, ffd7919·3096b26·c326d45); 탈출 청크가 버려지던 것 등(F25, 8bfe694).
  \item \textbf{그래서(결정):} 정본 반영: §84 보충 2 구현 기록(Task 21), §84 보충 8 구현 기록(Task 17: 명령 속도에 a를 곱한 뒤 가속·속도 상한, a = 0이면 Astra 수정 효과 0), §90 구현 기록 1--3, §91 구현 기록(우선순위 판정, 늦은 답은 늘 폐기, 예상 상태 = 지금 손끝 + 실행 청크 변위), §4·§58·§64 보충(C8). §6 결정 투영은 구현하지 않음.
\end{itemize}

\paragraph{C6. 무료 결합 사전 실행(셰이크다운) --- 배관은 되고 VLA가 바닥}
{\small 09-26 15:06 -- 16:28 KST}
\begin{itemize}
  \item \textbf{해 봄:} 사전 등록 9322853(15:06, P3 판정 없음, 0원): VLA = E-SR1c C1, Astra 대역 = Qwen3-VL-8B + \texttt{astra-\allowbreak{}couple@v1}(Astra 비슷한 지연 + 주입 시간 초과), DEV 0--5 $\times$ \{standard, dr\} $\times$ \{VLA 단독, 결합\}, 메인 GPU 1. 변경 1(f9903d7): 분석 지표만 추가.
  \item \textbf{결과:} bd3a8ba(16:28): 24/24편 충돌 0, 동시 1개 위반 0, 시간 초과·늦은 답 5(주입 6), stale 0, 스키마 오류 3/55, 편향 속도 최대 0.0084 $\leq$ 0.08 m/s, 가속 0.32 = 상한, 장부 0원. 성공 VLA 단독 0/12·결합 0/12.
  \item \textbf{결과:} 버그: B1 그리퍼 닫기를 풀 때 한 틱 14 mm 계단; B2 \texttt{max\_\allowbreak{}inflight} 기록이 증거가 못 됨; B3 편향·관문 로그 누락; B4 권한 a·§91·§90·움직임 줄이 이 판에 없음(a\_priv = 0에서 편향 0.42 mm 적용); B5(차단) VLA 단독 0/12 바닥 --- 짝 차가 정의상 0; B6 \texttt{b\_\allowbreak{}contradict} 사건 892회로 선택 단계 쉼이 안 켜짐; B7 편 끝 미응답 요청의 예약 청구; B8 Qwen 스키마 오류 5.5 \%·\texttt{grasped} 과대 주장 15.
  \item \textbf{그래서(결정):} B1--B3는 b0ee916(Task 17), B6·B7은 Task 21, B4는 Task 12·17--19로. B5 $\rightarrow$ 유료 E-Couple 전에 VLA 단독 폐루프 온전성 관문을 둔다.
\end{itemize}

\paragraph{C7. VLA 단독 0/12 분해 $\rightarrow$ E-VLA-solo: 의미 없음, 느리게 해도 안 됨, 교착 확인}
{\small 09-26 16:3x -- 21:22 KST}
\begin{itemize}
  \item \textbf{사용자} (ul 107, 16:3x KST (통제자 경유)): \intent{Vla는 혼자서 의미가 있는지도 보긴해야해 그리고 로본의 동작 속도자체도 조금 느리게 가져가도 좋을듯}
  \item \textbf{결과:} 진단(d752509, 18:44, 기존 기록만·CPU): 12편 중 8편이 런타임 교착 --- T1 모순 $\rightarrow$ M4 hold $\rightarrow$ 결정 비움 $\rightarrow$ 청크 버림 $\rightarrow$ 같은 모순. 방아쇠는 approach 중 청크 그리퍼 폭 77--80 mm가 열림 문턱 80.5 mm 바로 아래(3편), 들기 뒤 쥠 판정이 풀림(5편). 나머지 4편은 VLA 자체 실패(머그를 쳐서 탁자 밖 2, 머그 위에서 안 내려감 2). approach 중 모델 방향 답이 머그 쪽인 비율 36 \%(우연 수준).
  \item \textbf{해 봄:} E-VLA-solo 사전 등록(d752509, 18:44, 무료): E-SR1c C1, 7팔 $\times$ DEV 0--5 $\times$ standard·dr = 84편 --- C5·C3, 재생 1.0·0.75·0.5배, 계획기 규칙·제자리 기준선.
  \item \textbf{결과:} 21a5284(21:22): VLA 다섯 팔 모두 0/12, 계획기 규칙 12/12, 제자리 0/12. 지나침 87.7 $\rightarrow$ 52.9 $\rightarrow$ 17.9 mm(1.0·0.75·0.5배). C5 영구 hold 8/12 대 C3 0/12. C3는 들기 이상 10--11/12까지 가지만 놓기 성공 0. C3 관절 명령 계단(> 0.04 rad/틱) 103--135틱.
  \item \textbf{그래서(결정):} Q1 NOT\_MEANINGFUL(이 체크포인트의 VLA 단독은 성공을 못 냄) · Q2 NOT\_ADOPTED(재생 1.0배 유지, 정본 §92 보충 2 `속도 낮추기는 검증 후 채택') · Q3 CONFIRMED(교착 수정 필수). P3 짝 비교 보류, 유료 E-Couple 금지 유지.
\end{itemize}

\paragraph{C8. 교착 수정과 막힌 파드 검증}
{\small 09-26 21:57 -- 09-27 KST}
\begin{itemize}
  \item \textbf{해 봄:} Task 25(c835476, 21:57): T1 열림 히스테리시스(들어가기 80.5 mm, 나가기 75 mm), hold 3.0 s 뒤 한 스텝 탈출, hold 중 융합 청크 재생, VLA가 쥔 물체의 lift·carry 쥠 해제 $\rightarrow$ object\_lost 실패 경로, 융합 청크 전환 팔 틱당 0.04 rad 클램프. 수정 F25(8bfe694, 22:22): 탈출 청크 재생, 탈출 3연속 $\rightarrow$ \texttt{hold\_\allowbreak{}stuck} 실패, 힘만 떨어지면 \texttt{grip\_\allowbreak{}force\_\allowbreak{}low} 기록만(패드가 50.1 mm 아래로 닫혀야 object\_lost). 최종 검토 fddefd4.
  \item \textbf{진행 중·대기:} \tentative{파드 재측정(Task 25 C)과 결합 사전 실행 r2(Task 24)는 막혀 있다: HEAD의 \texttt{ser-\allowbreak{}A-\allowbreak{}min-\allowbreak{}3}이 기존 체크포인트를 모두 거부(serializer\_ok·last\_step\_ok·files\_ok). [결정 필요] ser-A-min-3 재학습 vs 옛 판 이식 vs 배선만 확인(handoff §0) --- 결합 학습 계획의 추천은 VLA-sim 재학습 먼저(C10 D2).}
\end{itemize}

\paragraph{C9. 상위$\leftrightarrow$VLA 의도 공유 프로토콜은 필수 핵심}
{\small 09-27 01:2x -- 01:30 KST}
\begin{itemize}
  \item \textbf{사용자} (ul 125, 01:2x KST): \intent{저 VLA는 0.33초 마다 스스로 움직여서 공백을 매우지만 상위에서 준 나는 이움직임동안 이러이런걸할거야가 같이 내려오면 그걸 기준으로 실행을 하는거야. 뭔소리인지 이해했지? 행동을 상위가 내 뱉을 때, 나는 이런이런 행동으로 할 생각이야 나는 물건을 향해 다다가는거야 하면서 행동 뱉고 그럼 그사이를 vla가 조종하고 그다음에 이번엔 잡을거야 하면 상위는 거기까지가고 vla는 그걸보고 이번엔 잡는구나 그래 다다가자 근데 언제 잡지?를 vla가하고 상위는 그걸보고 잡았구나 안잡았구나 이런 판단을 더 정확히 할 수 있으니까 그걸 하고. 이해했지? 이건 필수야 꼭 중요한 핵심이야}
  \item \textbf{그래서(결정):} 정본 §96 보충 2(ce7ba33, 01:30): 상위가 내는 모든 행동은 그 구간의 의도(§90 now/do/next)와 함께 내려온다; 답 사이(0.33 s)는 VLA가 그 의도를 기준으로 공백을 메운다; 쥐는 `언제'는 VLA 소유(§87); 상위는 VLA가 실제 한 일을 보고 잡았다/못 잡았다를 더 정확히 검증(§91·T1·V1h).
  \item \textbf{그래서(결정):} 학습 데이터에 대한 결과: 상위(35B)의 학습·추론 출력은 의도/구간 필드와 검증 판정을 반드시 포함 --- \texttt{astra-\allowbreak{}couple@v2} 같은 결합 형식. \texttt{astra-\allowbreak{}solo} 형식 실험(E-TEACH-L8, E-STRIP8, E-PT/ND)은 인식·조종 능력만 잰다.
\end{itemize}

\paragraph{C10. 결합 학습 계획 --- 엮을 부품은 있고, 모자란 것은 학습}
{\small 09-27 01:3x -- 01:39 KST}
\begin{itemize}
  \item \textbf{사용자} (ul 127, 01:3x KST): \intent{근데 이걸하려고 이미 과거에도 좀 정리를 했던게 있었을거야. 얼마나 그걸 잘 엮을 수 있을지 학습은 어떻게 할지가 이제 좀 제대로 잡혔지?}
  \item \textbf{결과:} \texttt{docs/\allowbreak{}research/\allowbreak{}coupled\_\allowbreak{}training\_\allowbreak{}plan\_\allowbreak{}2026-\allowbreak{}09-\allowbreak{}27.\allowbreak{}md}(a54170f, 01:39, 계획만·0원·GPU 0): 런타임 배선은 Task 1--26으로 끝. 없는 것 셋 --- (1) \texttt{ser-\allowbreak{}A-\allowbreak{}min-\allowbreak{}3} 형식으로 학습된 VLA 체크포인트 0개, (2) 결합 형식 상위 학습 데이터·참값 라벨러 \texttt{couple-\allowbreak{}truth} 없음(L8·E-PT·E-STRIP8은 모두 solo 형식), (3) VLA 끝-끝 학습·결합 폐루프 검증 없음. 어긋남 G1(상위 구간 9개 $\rightarrow$ VLA 줄 4개, 의도된 설계)·G2(결합 형식엔 공간 목표 없음)·G3(명시적 검증 필드 없음)·G4(상위 요청에 T1 판정 글 없음)·G5(VLA 학습에 늦은·틀린 의도 없음)·G6(체크포인트 0)·G7(§86 상수는 v1에서 잰 값)·G8(E-AT \texttt{quality:} 줄 없음).
  \item \textbf{그래서(결정):} 정본 §96 보충 3·4: 무료 사다리 V0 VLA 재학습 관문 $\rightarrow$ V1 VLA 단독 $\geq$ 4/12(영구 hold $\leq$ 1/12) $\rightarrow$ V2 결합 사전 실행 r2 $\rightarrow$ V3 E-TEACH-L8C 오프라인(\texttt{upper-\allowbreak{}couple@v1} = v2 + verify) $\rightarrow$ V4 결합 폐루프(학생 8B + VLA 대 8B 단독 동기 대 VLA 단독, DEV 40 짝) $\rightarrow$ V5 35B; V6 유료 E-Couple은 사용자 승인 뒤만. 자리 규칙 §88·§89 그대로. VLA 쪽 E2E는 D12.
  \item \textbf{진행 중·대기:} \tentative{사용자 결정 3개(추천): D1 학생 형식에 \texttt{verify} 필드 추가(예), D2 GPU 순서 --- ser-A-min-3 VLA-sim 재학습을 가장 먼저(예), D3 VLA 구간 줄은 지금 4단계 유지하고 V4 쥐기 시점 오차로 다시 판단. book 01의 VLA-R3·E-VLA-E2E·E-TEACH-L8C·E-COUPLE-8B·E-TEACH-35C는 모두 [예정].}
\end{itemize}
